\documentclass[aspectratio=169,10pt]{ctexbeamer}
\usepackage{amsmath,amssymb,mathtools}
\usepackage{booktabs}
\usepackage{xcolor}
\usepackage{hyperref}
\usepackage{graphicx}

\usetheme{Madrid}
\setbeamertemplate{navigation symbols}{}

\definecolor{Main}{RGB}{22,70,120}
\definecolor{Warn}{RGB}{170,50,35}
\definecolor{Soft}{RGB}{245,248,252}

\newcommand{\e}{\mathrm e}
\newcommand{\R}{\mathbb R}
\newcommand{\Q}{\mathbb Q}
\newcommand{\auditnote}[1]{\par\medskip{\color{Warn}\bfseries 原答案疏漏/补充：}#1}
\newcommand{\ocrnote}[1]{\par\medskip{\color{Warn}\bfseries 原文校验说明：}#1}

% Unified Yingxing Beamer footer and visual style.
\newcommand{\huayinglogo}{assets/logo1.jpg}
\newcommand{\huayingslogan}{英行培优，成绩无忧；英行相伴，刮目相看}

\setbeamersize{text margin left=8mm,text margin right=8mm}
\setbeamerfont{frametitle}{size=\large}
\setbeamercolor{structure}{fg=blue!70!black}
\setbeamercolor{frametitle}{bg=blue!75!black,fg=white}
\setbeamercolor{block title}{bg=blue!10,fg=black}
\setbeamercolor{block body}{bg=blue!2,fg=black}
\setbeamercolor{author in head/foot}{bg=blue!8,fg=black}
\setbeamercolor{title in head/foot}{bg=blue!8,fg=black}

\setbeamertemplate{footline}{%
  \leavevmode%
  \hbox{%
    \begin{beamercolorbox}[wd=.12\paperwidth,ht=3.8ex,dp=1.2ex,center]{author in head/foot}%
      \includegraphics[height=2.9ex]{\huayinglogo}%
    \end{beamercolorbox}%
    \begin{beamercolorbox}[wd=.76\paperwidth,ht=3.8ex,dp=1.2ex,center]{title in head/foot}%
      \scriptsize \huayingslogan
    \end{beamercolorbox}%
    \begin{beamercolorbox}[wd=.12\paperwidth,ht=3.8ex,dp=1.2ex,right]{author in head/foot}%
      \scriptsize \insertframenumber/\inserttotalframenumber\hspace*{1em}
    \end{beamercolorbox}%
  }%
}

\title{导数专题讲义}
\subtitle{对数、指数与多极值点函数最值问题}
\author{}
\date{}

\begin{document}

\begin{frame}
  \titlepage
\end{frame}

\begin{frame}{审计说明}
\small
本版本按两份原 PDF 的顺序重新整理：
\begin{itemize}
  \item 保留所有方法总结、重要模型、例题、巩固练习和参考答案思路。
  \item 每道例题或习题均先给出完整题目，随后紧跟答案解释页。
  \item 若原答案中有“略”、公式缺失、OCR 乱码、推理跳步或重复题，均在答案页标明。
  \item 页码顺序以原 PDF 出现顺序为准；第二专题原文例题编号从例 3 跳到例 6，本文保留该编号并说明。
\end{itemize}
\end{frame}

\begin{frame}{目录}
\tableofcontents
\end{frame}

\section{专题一：对数、指数与指对分离}

\subsection{考点一：对数单身狗}

\begin{frame}[allowframebreaks]{方法总结：对数单身狗}
\small
处理含对数函数的不等式或方程时，若直接对
\[
f(x)\ln x+g(x)
\]
反复求导，导数中仍含有 $\ln x$，计算通常很繁。核心做法是把对数项“单独分离”：
\[
f(x)\ln x+g(x)=f(x)\left(\ln x+\frac{g(x)}{f(x)}\right).
\]
这样研究
\[
\ln x+\frac{g(x)}{f(x)}
\]
时，其导函数为
\[
\frac1x+\left(\frac{g(x)}{f(x)}\right)',
\]
通常不再含超越函数，极值点和单调性更容易判断。

\framebreak

两种常见形式：
\[
f(x)>0,\quad f(x)\ln x+g(x)>0
\Longleftrightarrow
\ln x+\frac{g(x)}{f(x)}>0.
\]
\[
f(x)\ne0,\quad f(x)\ln x+g(x)=0
\Longleftrightarrow
\ln x+\frac{g(x)}{f(x)}=0.
\]
这就是原讲义中的“对数单身狗”：把 $\ln x$ 留出来，研究剩余部分。
\end{frame}

\begin{frame}{例 1 题目}
\small
（2016 全国 II）已知函数
\[
f(x)=(x+1)\ln x-a(x-1),\qquad x>0.
\]
\begin{enumerate}
  \item 当 $a=4$ 时，求曲线 $y=f(x)$ 在 $(1,f(1))$ 处的切线方程；
  \item 若当 $x\in(1,+\infty)$ 时，$f(x)>0$，求 $a$ 的取值范围。
\end{enumerate}
\end{frame}

\begin{frame}[allowframebreaks]{例 1 答案解释}
\small
当 $a=4$ 时，
\[
f(1)=0,\qquad f'(x)=\ln x+\frac1x+1-a,
\]
故
\[
f'(1)=2-4=-2.
\]
切线为
\[
y-0=-2(x-1),
\]
即
\[
\boxed{2x+y-2=0}.
\]

\framebreak

第二问中，$x>1$ 时
\[
f(x)>0
\Longleftrightarrow
\ln x-\frac{a(x-1)}{x+1}>0.
\]
令
\[
g(x)=\ln x-\frac{a(x-1)}{x+1},
\]
则
\[
g'(x)=\frac1x-\frac{2a}{(x+1)^2}
=\frac{x^2+2(1-a)x+1}{x(x+1)^2}.
\]
且 $g(1)=0$。

\framebreak

若 $a\le2$，则
\[
x^2+2(1-a)x+1\ge x^2-2x+1=(x-1)^2\ge0,
\]
所以 $g'(x)\ge0$，$g(x)>g(1)=0$ 对 $x>1$ 成立。

若 $a>2$，方程
\[
x^2+2(1-a)x+1=0
\]
有两根
\[
x_1=a-1-\sqrt{(a-1)^2-1},\qquad
x_2=a-1+\sqrt{(a-1)^2-1},
\]
且 $x_1x_2=1,\ x_2>1$。因此 $g$ 在 $(1,x_2)$ 上递减，故 $g(x)<g(1)=0$，不满足题意。

\[
\boxed{a\le2}.
\]
\end{frame}

\begin{frame}{例 2 题目}
\small
已知函数
\[
f(x)=\frac{a\ln x}{x+1}+\frac bx,\qquad x>0,
\]
曲线 $y=f(x)$ 在点 $(1,f(1))$ 处的切线方程为
\[
x+2y-3=0.
\]
\begin{enumerate}
  \item 求 $a,b$ 的值；
  \item 证明：当 $x>0$ 且 $x\ne1$ 时，
  \[
  f(x)>\frac{\ln x}{x-1}.
  \]
\end{enumerate}
\end{frame}

\begin{frame}[allowframebreaks]{例 2 答案解释}
\small
由
\[
f(1)=b,\qquad f'(x)=
\frac{a\left(\frac{x+1}{x}-\ln x\right)}{(x+1)^2}-\frac b{x^2},
\]
切线 $x+2y-3=0$ 斜率为 $-\frac12$，且过 $(1,1)$，故
\[
f(1)=1,\qquad f'(1)=-\frac12.
\]
于是
\[
b=1,\qquad \frac a2-b=-\frac12,
\]
解得
\[
\boxed{a=1,\quad b=1}.
\]

\framebreak

由第一问
\[
f(x)=\frac{\ln x}{x+1}+\frac1x.
\]
计算差值：
\[
f(x)-\frac{\ln x}{x-1}
=\frac1{1-x^2}\left(2\ln x-\frac{x^2-1}{x}\right).
\]
令
\[
h(x)=2\ln x-\frac{x^2-1}{x}=2\ln x-x+\frac1x.
\]
则
\[
h'(x)=\frac2x-1-\frac1{x^2}
=-\frac{(x-1)^2}{x^2}\le0.
\]
因此 $h$ 在 $(0,+\infty)$ 上单调递减，且 $h(1)=0$。

当 $0<x<1$ 时，$h(x)>0$ 且 $1-x^2>0$；当 $x>1$ 时，$h(x)<0$ 且 $1-x^2<0$。两种情况下差值均为正。
\[
\boxed{f(x)>\frac{\ln x}{x-1}}.
\]
\end{frame}

\begin{frame}{巩固 1 题目}
\small
若不等式
\[
x\ln x\ge a(x-1)
\]
对所有 $x\ge1$ 都成立，求实数 $a$ 的取值范围。
\end{frame}

\begin{frame}[allowframebreaks]{巩固 1 答案解释}
\small
原不等式等价于
\[
\ln x-\frac{a(x-1)}x\ge0,\qquad x\ge1.
\]
令
\[
h(x)=\ln x-\frac{a(x-1)}x.
\]
则
\[
h'(x)=\frac1x-\frac a{x^2}=\frac{x-a}{x^2},\qquad h(1)=0.
\]
若 $a\le1$，则 $x-a\ge0$ 对 $x\ge1$ 恒成立，故 $h$ 在 $[1,+\infty)$ 上递增，$h(x)\ge h(1)=0$。

若 $a>1$，则 $h$ 在 $(1,a)$ 上递减，因而 $h(x)<h(1)=0$，不合题意。
\[
\boxed{a\le1}.
\]
\end{frame}

\begin{frame}{巩固 2 题目}
\small
（2017 全国 II）已知函数
\[
f(x)=ax^2-ax-x\ln x,\qquad x>0,
\]
且 $f(x)\ge0$。
\begin{enumerate}
  \item 求 $a$；
  \item 证明：$f(x)$ 存在唯一的极大值点 $x_0$，且
  \[
  e^{-2}<f(x_0)<2^{-2}.
  \]
\end{enumerate}
\end{frame}

\begin{frame}[allowframebreaks]{巩固 2 答案解释}
\small
写成
\[
f(x)=x(ax-a-\ln x)=xg(x).
\]
由于 $x>0$，$f(x)\ge0$ 等价于 $g(x)\ge0$。又
\[
g(1)=0,\qquad g'(x)=a-\frac1x,\qquad g'(1)=a-1.
\]
若 $g(x)\ge0$ 且 $g(1)=0$，则 $x=1$ 是 $g$ 的最小点，故 $g'(1)=0$，得
\[
\boxed{a=1}.
\]
此时
\[
g'(x)=1-\frac1x,
\]
所以 $g$ 在 $(0,1)$ 上递减，在 $(1,+\infty)$ 上递增，$g(x)\ge0$ 成立。

\framebreak

当 $a=1$ 时
\[
f(x)=x^2-x-x\ln x,\qquad f'(x)=2x-2-\ln x.
\]
令
\[
h(x)=2x-2-\ln x.
\]
则
\[
h'(x)=2-\frac1x.
\]
在 $(0,\frac12)$ 上 $h'<0$，在 $(\frac12,+\infty)$ 上 $h'>0$。又
\[
h(e^{-2})=2e^{-2}>0,\quad h\!\left(\frac12\right)<0,\quad h(1)=0.
\]
故 $h$ 在 $(0,\frac12)$ 有唯一零点 $x_0$，在 $[\frac12,+\infty)$ 有唯一零点 $1$。

\framebreak

由符号可知，$f'(x)$ 在 $(0,x_0)$ 为正，在 $(x_0,1)$ 为负，在 $(1,+\infty)$ 为正，因此 $x_0$ 是唯一极大值点。

由 $f'(x_0)=0$ 得
\[
\ln x_0=2x_0-2.
\]
所以
\[
f(x_0)=x_0^2-x_0-x_0\ln x_0
=x_0^2-x_0-x_0(2x_0-2)
=x_0(1-x_0).
\]
因为 $x_0\in(0,1)$，所以 $f(x_0)<\frac14=2^{-2}$。

又 $e^{-1}\in(0,1)$ 且不是驻点，原函数在 $(0,1)$ 的最大值取于 $x_0$，故
\[
f(x_0)>f(e^{-1})=e^{-2}.
\]
综上
\[
\boxed{e^{-2}<f(x_0)<2^{-2}}.
\]
\end{frame}

\begin{frame}{巩固 3 题目}
\small
（2018 全国 III）已知函数
\[
f(x)=(2+x+ax^2)\ln(1+x)-2x,\qquad x>-1.
\]
\begin{enumerate}
  \item 若 $a=0$，证明：当 $-1<x<0$ 时，$f(x)<0$；当 $x>0$ 时，$f(x)>0$；
  \item 若 $x=0$ 是 $f(x)$ 的极大值点，求 $a$。
\end{enumerate}
\end{frame}

\begin{frame}[allowframebreaks]{巩固 3 答案解释}
\small
当 $a=0$ 时，
\[
f(x)=(2+x)\ln(1+x)-2x.
\]
求导得
\[
f'(x)=\ln(1+x)-\frac{x}{1+x}.
\]
令
\[
g(x)=\ln(1+x)-\frac{x}{1+x},
\]
则
\[
g'(x)=\frac{x}{(1+x)^2}.
\]
故 $g$ 在 $(-1,0)$ 上递减，在 $(0,+\infty)$ 上递增，且 $g(0)=0$，于是 $g(x)\ge0$，等号仅在 $x=0$ 取到。由 $f(0)=0$ 可知 $f$ 在 $(-1,+\infty)$ 上递增，因此
\[
-1<x<0\Rightarrow f(x)<0,\qquad x>0\Rightarrow f(x)>0.
\]

\framebreak

第二问中，若 $a\ge0$，由第一问结论可知 $x>0$ 时
\[
f(x)\ge (2+x)\ln(1+x)-2x>0=f(0),
\]
这与 $x=0$ 是极大值点矛盾，所以必须 $a<0$。

令
\[
h(x)=\frac{f(x)}{2+x+ax^2}
=\ln(1+x)-\frac{2x}{2+x+ax^2},
\]
在 $x=0$ 附近分母为正，故 $h$ 与 $f$ 局部同号，且 $h(0)=0$。要使 $x=0$ 为极大值点，只需 $h$ 在 $0$ 附近左增右减。

计算
\[
h'(x)=\frac{x^2(a^2x^2+4ax+6a+1)}{(x+1)(ax^2+x+2)^2}.
\]
因此 $x=0$ 附近 $h'(x)$ 的符号取决于 $6a+1$。只有
\[
6a+1=0
\]
时，才可出现左增右减。代入可得
\[
h'(x)=\frac{x^3(x-24)}{(x+1)(x^2-6x-12)^2},
\]
从而 $(-1,0)$ 上 $h'>0$，$(0,1)$ 上 $h'<0$，故 $x=0$ 确为极大值点。
\[
\boxed{a=-\frac16}.
\]
\end{frame}

\subsection{考点二：指数找基友}

\begin{frame}{方法总结：指数找基友}
\small
处理含指数函数的不等式时，常把指数型函数与多项式函数“结合”：
\[
\e^x+f(x)>0
\Longleftrightarrow
1+\frac{f(x)}{\e^x}>0.
\]
此时
\[
\left(1+\frac{f(x)}{\e^x}\right)'
=\frac{f'(x)-f(x)}{\e^x},
\]
若 $f$ 是多项式，导数结构会明显简化。

同理，
\[
\e^x+f(x)=0
\Longleftrightarrow
1+\frac{f(x)}{\e^x}=0.
\]
这就是“指数找基友”：让指数项和代数项配对，把求导后的函数降为更可控的代数形式。
\end{frame}

\begin{frame}{例 3 题目}
\small
（2018 全国 II）已知函数
\[
f(x)=\e^x-ax^2.
\]
\begin{enumerate}
  \item 若 $a=1$，证明：当 $x\ge0$ 时，$f(x)\ge1$；
  \item 若 $f(x)$ 在 $(0,+\infty)$ 只有一个零点，求 $a$。
\end{enumerate}
\end{frame}

\begin{frame}[allowframebreaks]{例 3 答案解释}
\small
当 $a=1$ 时，要证
\[
\e^x-x^2\ge1
\]
等价于
\[
(x^2+1)\e^{-x}\le1.
\]
令
\[
g(x)=(x^2+1)\e^{-x}-1,
\]
则
\[
g'(x)=-(x-1)^2\e^{-x}\le0.
\]
故 $g$ 在 $[0,+\infty)$ 上单调递减，且 $g(0)=0$，所以 $g(x)\le0$，即
\[
\boxed{\e^x-x^2\ge1}.
\]

\framebreak

第二问可采用参变分离：
\[
\e^x-ax^2=0
\Longleftrightarrow
a=\frac{\e^x}{x^2}=:\varphi(x),\qquad x>0.
\]
求导：
\[
\varphi'(x)=\frac{\e^x(x-2)}{x^3}.
\]
故 $\varphi$ 在 $(0,2)$ 上递减，在 $(2,+\infty)$ 上递增，最小值为
\[
\varphi(2)=\frac{\e^2}{4}.
\]
当 $a=\frac{\e^2}{4}$ 时，方程有唯一正根；当 $a<\frac{\e^2}{4}$ 无正根，当 $a>\frac{\e^2}{4}$ 有两个正根。
\[
\boxed{a=\frac{\e^2}{4}}.
\]
\end{frame}

\begin{frame}{例 4 题目}
\small
（2020 全国 I）已知函数
\[
f(x)=\e^x+ax^2-x.
\]
\begin{enumerate}
  \item 当 $a=1$ 时，讨论 $f(x)$ 的单调性；
  \item 当 $x\ge0$ 时，
  \[
  f(x)\ge\frac12x^3+1,
  \]
  求 $a$ 的取值范围。
\end{enumerate}
\ocrnote{原图中不等号易看成严格大于；若写成 $>$，在 $x=0$ 时两边同为 $1$，题目不成立。按解析和高考原题应为 $\ge$。}
\end{frame}

\begin{frame}[allowframebreaks]{例 4 答案解释}
\small
当 $a=1$ 时
\[
f(x)=\e^x+x^2-x,\qquad f'(x)=\e^x+2x-1.
\]
又
\[
f''(x)=\e^x+2>0,
\]
故 $f'$ 严格递增，且 $f'(0)=0$。于是
\[
x<0\Rightarrow f'(x)<0,\qquad x>0\Rightarrow f'(x)>0.
\]
所以 $f$ 在 $(-\infty,0)$ 上单调递减，在 $(0,+\infty)$ 上单调递增。

\framebreak

第二问等价于
\[
\e^x+ax^2-x\ge\frac12x^3+1,\qquad x\ge0.
\]
当 $x=0$ 时恒成立。若 $x>0$，分离参数得
\[
a\ge-\frac{\e^x-\frac12x^3-x-1}{x^2}.
\]
令
\[
g(x)=-\frac{\e^x-\frac12x^3-x-1}{x^2}.
\]
利用 $\e^x-\frac12x^2-x-1\ge0$ 可得
\[
g'(x)=-\frac{(x-2)\left(\e^x-\frac12x^2-x-1\right)}{x^3}.
\]
因此 $g$ 在 $(0,2)$ 上递增，在 $(2,+\infty)$ 上递减，最大值为
\[
g(2)=\frac{7-\e^2}{4}.
\]
于是
\[
\boxed{a\ge\frac{7-\e^2}{4}}.
\]
\auditnote{原答案中有一处把 $f''(x)=\e^x+2>0$ 误写成 $f'(x)=\e^x+2>0$，本版已修正。}
\end{frame}

\begin{frame}{巩固 4 题目}
\small
已知函数
\[
f(x)=\e^x-1-x-ax^2.
\]
当 $x\ge0$ 时，$f(x)\ge0$ 恒成立，求实数 $a$ 的取值范围。
\end{frame}

\begin{frame}[allowframebreaks]{巩固 4 答案解释}
\small
若
\[
f(x)=\e^x-1-x-ax^2\ge0,\qquad x\ge0,
\]
则当 $x>0$ 时
\[
a\le\frac{\e^x-1-x}{x^2}.
\]
记
\[
\psi(x)=\frac{\e^x-1-x}{x^2}.
\]
由泰勒展开可知
\[
\lim_{x\to0^+}\psi(x)=\frac12.
\]
也可按原答案使用 $\e^x\ge x+1$：
当 $a\le\frac12$ 时，
\[
f'(x)=\e^x-1-2ax\ge x-2ax=(1-2a)x\ge0,
\]
且 $f(0)=0$，故 $f(x)\ge0$。

若 $a>\frac12$，则在 $x>0$ 很小时，$f'(x)<0$ 或等价地
\[
\e^x-1-x-ax^2=\left(\frac12-a\right)x^2+o(x^2)<0,
\]
不合题意。
\[
\boxed{a\le\frac12}.
\]
\end{frame}

\begin{frame}{巩固 5 题目}
\small
已知函数
\[
f(x)=\e^{-x}+ax,\qquad a\in\R.
\]
\begin{enumerate}
  \item 讨论 $f(x)$ 的最值；
  \item 若 $a=0$，证明：
  \[
  f(x)>-\frac12x^2+\frac58.
  \]
\end{enumerate}
\end{frame}

\begin{frame}[allowframebreaks]{巩固 5 答案解释}
\small
有
\[
f'(x)=-\e^{-x}+a.
\]
若 $a\le0$，则 $f'(x)<0$，$f$ 在 $\R$ 上单调递减，不存在最大值和最小值。

若 $a>0$，由 $f'(x)=0$ 得
\[
x=-\ln a.
\]
此时 $f$ 在 $(-\infty,-\ln a)$ 上递减，在 $(-\ln a,+\infty)$ 上递增，故最小值为
\[
f(-\ln a)=a-a\ln a,
\]
无最大值。

\framebreak

当 $a=0$ 时，$f(x)=\e^{-x}$。要证
\[
\e^{-x}>-\frac12x^2+\frac58
\]
等价于
\[
(5-4x^2)\e^x<8.
\]
设
\[
h(x)=(5-4x^2)\e^x.
\]
若 $5-4x^2\le0$，结论显然。若 $-\frac{\sqrt5}{2}<x<\frac{\sqrt5}{2}$，
\[
h'(x)=(-4x^2-8x+5)\e^x=-(2x-1)(2x+5)\e^x.
\]
所以 $h$ 在 $(-\frac{\sqrt5}{2},\frac12)$ 上递增，在 $(\frac12,\frac{\sqrt5}{2})$ 上递减，最大值
\[
h\left(\frac12\right)=4\sqrt{\e}<8.
\]
故不等式成立。
\end{frame}

\begin{frame}{巩固 6 题目}
\small
已知函数
\[
f(x)=a(x-1),\qquad g(x)=(ax-1)\e^x,\qquad a\in\R.
\]
\begin{enumerate}
  \item 证明：存在唯一实数 $a$，使得直线 $y=f(x)$ 和曲线 $y=g(x)$ 相切；
  \item 若不等式 $f(x)>g(x)$ 有且只有两个整数解，求 $a$ 的取值范围。
\end{enumerate}
\end{frame}

\begin{frame}[allowframebreaks]{巩固 6 答案解释}
\small
有
\[
f'(x)=a,\qquad g'(x)=(ax+a-1)\e^x.
\]
设切点为 $(x_0,y_0)$，则
\[
a(x_0-1)=(ax_0-1)\e^{x_0},
\]
即
\[
a(x_0\e^{x_0}-x_0+1)=\e^{x_0}. \tag{1}
\]
又因相切，
\[
a=(ax_0+a-1)\e^{x_0},
\]
整理得
\[
a(x_0\e^{x_0}+\e^{x_0}-1)=\e^{x_0}. \tag{2}
\]
联立 (1)(2) 得
\[
\e^{x_0}+x_0-2=0.
\]
令 $H(x)=\e^x+x-2$，则 $H'(x)=\e^x+1>0$，且 $H(0)<0,H(1)>0$，故唯一 $x_0\in(0,1)$，继而唯一 $a$。

\framebreak

第二问中
\[
f(x)>g(x)
\Longleftrightarrow
a(x-1)>(ax-1)\e^x.
\]
等价于
\[
a\left(x-\frac{x-1}{\e^x}\right)<1.
\]
令
\[
m(x)=x-\frac{x-1}{\e^x}.
\]
则
\[
m'(x)=\frac{\e^x+x-2}{\e^x}.
\]
由第一问可知 $m$ 在 $(-\infty,x_0)$ 递减，在 $(x_0,+\infty)$ 递增，且 $x_0\in(0,1)$。

当 $x\le0$ 或 $x\ge1$ 且 $x\in\mathbb Z$ 时，$m(x)\ge1$。要使整数解恰为 $0,1$，需
\[
m(2)\ge\frac1a,\qquad m(-1)\ge\frac1a,
\]
解得
\[
\boxed{a\in\left[\frac{\e^2}{2\e^2-1},1\right)}.
\]
\auditnote{原答案末尾给成 $\left[\frac{\e^2}{2\e^2-1},+\infty\right)$，但前面分类已排除 $a\ge1$，否则没有符合要求的两个整数解。正确范围应为 $\left[\frac{\e^2}{2\e^2-1},1\right)$。}
\end{frame}

\subsection{考点三：指对在一起，常常要分手}

\begin{frame}{方法总结：指对分离}
\small
当式子中同时出现 $\e^x$ 与 $\ln x$ 时，直接求导往往复杂。原讲义建议将指数和对数拆开，例如
\[
(\e^x\ln x-f(x))'=\e^x\ln x+\frac{\e^x}{x}-f'(x).
\]
若求导后仍有 $\e^x,\ln x$ 混在一起，可构造
\[
\frac{\e^x\ln x-f(x)}{\e^x}
=\ln x-\frac{f(x)}{\e^x},
\]
从而把指数部分先消去，再单独研究对数部分。
\end{frame}

\begin{frame}{例 5 题目}
\small
（2014 全国 I）设函数
\[
f(x)=a\e^x\ln x+\frac{b\e^{x-1}}x,\qquad x>0.
\]
曲线 $y=f(x)$ 在点 $(1,f(1))$ 处的切线为
\[
y=\e(x-1)+2.
\]
\begin{enumerate}
  \item 求 $a,b$；
  \item 证明：$f(x)>1$。
\end{enumerate}
\end{frame}

\begin{frame}[allowframebreaks]{例 5 答案解释}
\small
求导：
\[
f'(x)=a\e^x\left(\ln x+\frac1x\right)+\frac{b\e^{x-1}(x-1)}{x^2}.
\]
切线过 $(1,2)$ 且斜率为 $\e$，所以
\[
f(1)=2,\qquad f'(1)=\e.
\]
代入得
\[
b=2,\qquad a\e=\e,
\]
故
\[
\boxed{a=1,\quad b=2}.
\]

\framebreak

此时
\[
f(x)=\e^x\ln x+\frac{2\e^{x-1}}x.
\]
证明 $f(x)>1$ 等价于
\[
x\ln x>x\e^{-x}-\frac2\e.
\]
令
\[
G(x)=x\ln x,\qquad H(x)=x\e^{-x}-\frac2\e.
\]
则
\[
G'(x)=1+\ln x,
\]
所以 $G$ 在 $(0,\frac1\e)$ 递减，在 $(\frac1\e,+\infty)$ 递增，最小值为
\[
G\left(\frac1\e\right)=-\frac1\e.
\]
而
\[
H'(x)=\e^{-x}(1-x),
\]
故 $H$ 在 $(0,1)$ 递增，在 $(1,+\infty)$ 递减，最大值为
\[
H(1)=-\frac1\e.
\]
两函数取到极值的位置不同，因此对所有 $x>0$ 有 $G(x)>H(x)$，从而
\[
\boxed{f(x)>1}.
\]
\end{frame}

\begin{frame}{例 6 题目}
\small
已知函数
\[
f(x)=\frac1x+a\ln x,\qquad g(x)=\frac{\e^x}{x},\qquad x>0.
\]
\begin{enumerate}
  \item 讨论函数 $f(x)$ 的单调性；
  \item 证明：当 $a=1$ 时，
  \[
  f(x)+g(x)-\left(1+\frac{\e}{x^2}\right)\ln x>\e.
  \]
\end{enumerate}
\end{frame}

\begin{frame}[allowframebreaks]{例 6 答案解释}
\small
有
\[
f'(x)=-\frac1{x^2}+\frac ax=\frac{ax-1}{x^2}.
\]
若 $a\le0$，则 $f'(x)<0$，$f$ 在 $(0,+\infty)$ 上单调递减。

若 $a>0$，则
\[
f'(x)<0\Longleftrightarrow 0<x<\frac1a,\qquad
f'(x)>0\Longleftrightarrow x>\frac1a.
\]
故 $f$ 在 $(0,\frac1a)$ 上递减，在 $(\frac1a,+\infty)$ 上递增。

\framebreak

当 $a=1$ 时，要证
\[
\frac1x+\frac{\e^x}{x}-\left(1+\frac{\e}{x^2}\right)\ln x>\e.
\]
等价于
\[
\e^x-\e x+1>\frac{\e\ln x}{x}.
\]
令
\[
F(x)=\e^x-\e x+1,\qquad F'(x)=\e^x-\e.
\]
则 $F$ 在 $(0,1)$ 上递减，在 $(1,+\infty)$ 上递增，最小值 $F(1)=1$。

令
\[
G(x)=\frac{\e\ln x}{x},\qquad G'(x)=\frac{\e(1-\ln x)}{x^2}.
\]
则 $G$ 在 $(0,\e)$ 上递增，在 $(\e,+\infty)$ 上递减，最大值 $G(\e)=1$。
由于 $F$ 与 $G$ 取极值点不同，故 $F(x)>G(x)$，原不等式成立。
\end{frame}

\begin{frame}{巩固 7 题目}
\small
设函数
\[
f(x)=\frac{\ln x+1}{x}.
\]
求证：当 $x>1$ 时，
\[
\frac{f(x)}{\e+1}>
\frac{2\e^{x-1}}{(x+1)(x\e^x+1)}.
\]
\end{frame}

\begin{frame}[allowframebreaks]{巩固 7 答案解释}
\small
不等式等价于
\[
\frac1{\e+1}\cdot\frac{(x+1)(\ln x+1)}x
>
\frac{2\e^{x-1}}{x\e^x+1}.
\]
分别构造
\[
G(x)=\frac{(x+1)(\ln x+1)}x,\qquad
H(x)=\frac{2\e^{x-1}}{x\e^x+1}.
\]
对 $G$，
\[
G'(x)=\frac{x-\ln x}{x^2}.
\]
令 $\varphi(x)=x-\ln x$，则
\[
\varphi'(x)=1-\frac1x=\frac{x-1}x>0\quad(x>1),
\]
故 $\varphi(x)>\varphi(1)=1>0$，于是 $G'(x)>0$，$G(x)>G(1)=2$。

对 $H$，
\[
H'(x)=\frac{2\e^{x-1}(1-\e^x)}{(x\e^x+1)^2}<0\quad(x>1),
\]
所以 $H(x)<H(1)=\frac2{\e+1}$。因此
\[
\frac{G(x)}{\e+1}>H(x),
\]
证毕。
\end{frame}

\begin{frame}{巩固 8 题目}
\small
已知函数
\[
f(x)=\e^x-a\ln x-a,\qquad a>0,\ x>0.
\]
\begin{enumerate}
  \item 当 $a>\e$ 时，求函数 $f(x)$ 的极值；
  \item 求证：
  \[
  \e^{2x-2}-\e^{x-1}\ln x-x\ge0.
  \]
\end{enumerate}
\ocrnote{原题干写作“当 $a>\e$ 时”，但原答案实际按 $a=\e$ 处理；本版先给出 $a>\e$ 的正确极值讨论，再说明第 (2) 问使用的是 $a=\e$ 时的辅助结论。}
\end{frame}

\begin{frame}[allowframebreaks]{巩固 8 答案解释}
\small
先处理题干中的 $a>\e$。有
\[
f'(x)=\e^x-\frac ax.
\]
且
\[
f''(x)=\e^x+\frac a{x^2}>0,
\]
故 $f'$ 严格递增。方程 $f'(x)=0$ 等价于
\[
x\e^x=a.
\]
因为 $x\e^x$ 在 $(0,+\infty)$ 上严格递增，所以存在唯一 $\alpha>1$ 满足
\[
\alpha\e^\alpha=a.
\]
于是 $f$ 在 $(0,\alpha)$ 上递减，在 $(\alpha,+\infty)$ 上递增，唯一极小值为
\[
f(\alpha)=\e^\alpha-a\ln\alpha-a,
\]
无极大值。

\auditnote{原答案把本问当成 $a=\e$，直接写 $x=1$ 为极小值点。若题干确为 $a>\e$，这个结论不正确。}

\framebreak

第 (2) 问原答案使用的是 $a=\e$ 时的辅助结论：
\[
\e^x-\e\ln x-\e\ge0,
\]
即
\[
\e^x-\e\ln x\ge\e.
\]
要证
\[
\e^{2x-2}-\e^{x-1}\ln x-x\ge0
\]
等价于
\[
\e^x-\e\ln x\ge \frac{x}{\e^{x-2}}.
\]
令
\[
g(x)=\frac{x}{\e^{x-2}}=x\e^{2-x}.
\]
则
\[
g'(x)=(1-x)\e^{2-x}.
\]
故 $g$ 在 $(0,1)$ 上递增，在 $(1,+\infty)$ 上递减，最大值 $g(1)=\e$。于是
\[
\e^x-\e\ln x\ge\e\ge g(x),
\]
结论成立。
\end{frame}

\begin{frame}{巩固 9 题目}
\small
已知函数
\[
f(x)=\frac1x+a\ln x,\qquad g(x)=\frac{\e^x}{x}.
\]
\begin{enumerate}
  \item 讨论函数 $f(x)$ 的单调性；
  \item 证明：当 $a=1$ 时，
  \[
  f(x)+g(x)-\left(1+\frac{\e}{x^2}\right)\ln x>\e.
  \]
\end{enumerate}
\ocrnote{原 PDF 此题作为考点三巩固练习第 3 题再次出现，与例 6 完全重复；本版本按原顺序保留。}
\end{frame}

\begin{frame}{巩固 9 答案解释}
\small
本题与例 6 重复，解法完全相同：
\[
f'(x)=\frac{ax-1}{x^2}.
\]
若 $a\le0$，$f$ 在 $(0,+\infty)$ 上递减；若 $a>0$，$f$ 在 $(0,\frac1a)$ 上递减，在 $(\frac1a,+\infty)$ 上递增。

当 $a=1$ 时，转化为证明
\[
\e^x-\e x+1>\frac{\e\ln x}{x}.
\]
左侧函数 $F(x)=\e^x-\e x+1$ 在 $x=1$ 取最小值 $1$，右侧函数 $G(x)=\frac{\e\ln x}{x}$ 在 $x=\e$ 取最大值 $1$，且取等位置不同，故严格不等式成立。
\end{frame}

\section{专题二：多极值点函数最值问题}

\begin{frame}[allowframebreaks]{基本原理：五种途径}
\small
处理多极值点问题时，关键是利用导函数零点之间的关系来降元。

\begin{enumerate}
  \item 消元，找到极值点 $x_1,x_2$ 的关系，常由导函数零点方程和韦达定理得到，再把目标量化为单个极值点的函数。
  \item 消元，构建极值点 $x_1,x_2$ 与参数 $a$ 的函数关系，最终转化为参数函数。
  \item 构造齐次式消元。若 $f(x)=\ln x-ax$ 且 $f(x_1)=f(x_2)=0$，设 $t=\frac{x_2}{x_1}>1$，可把 $x_1,x_2$ 表成 $t$ 的函数。
  \item 使用对数均值不等式：
  \[
  \sqrt{ab}\le L(a,b)=\frac{a-b}{\ln a-\ln b}\le\frac{a+b}{2}.
  \]
  \item 关注两个重要模型：
  \[
  f(x)=x-\frac1x-a\ln x,\qquad
  g(x)=(x+1)\ln x-a(x-1).
  \]
  若 $f(x)=0$，则 $f(1/x)=0$，常带来 $x_1x_3=1$ 的隐含关系。
\end{enumerate}
\end{frame}

\begin{frame}{例 1 题目}
\small
（2018 全国 I）已知函数
\[
f(x)=\frac1x-x+a\ln x,\qquad x>0.
\]
\begin{enumerate}
  \item 讨论 $f(x)$ 的单调性；
  \item 若 $f(x)$ 存在两个极值点 $x_1,x_2$，证明：
  \[
  \frac{f(x_1)-f(x_2)}{x_1-x_2}<a-2.
  \]
\end{enumerate}
\end{frame}

\begin{frame}[allowframebreaks]{例 1 答案解释}
\small
有
\[
f'(x)=-\frac1{x^2}-1+\frac ax
=-\frac{x^2-ax+1}{x^2}.
\]
若 $a\le2$，则 $x^2-ax+1\ge0$ 对 $x>0$ 成立，故 $f'(x)\le0$，$f$ 单调递减。

若 $a>2$，方程
\[
x^2-ax+1=0
\]
有两个正根 $x_1<x_2$，且 $x_1x_2=1$。此时 $f$ 在 $(0,x_1)$ 递减，$(x_1,x_2)$ 递增，$(x_2,+\infty)$ 递减。

\framebreak

当存在两个极值点时 $a>2$，且 $x_1x_2=1$。不妨设 $x_1<x_2$，则 $x_2>1$。
计算
\[
\frac{f(x_1)-f(x_2)}{x_1-x_2}
=-2+a\frac{\ln x_1-\ln x_2}{x_1-x_2}.
\]
利用 $x_1=1/x_2$，得
\[
\frac{\ln x_1-\ln x_2}{x_1-x_2}
=\frac{-2\ln x_2}{1/x_2-x_2}.
\]
要证原式 $<a-2$，等价于
\[
\frac1{x_2}-x_2+2\ln x_2<0.
\]
令
\[
\phi(x)=\frac1x-x+2\ln x.
\]
则
\[
\phi'(x)=-\frac1{x^2}-1+\frac2x=-\frac{(x-1)^2}{x^2}\le0.
\]
故 $\phi$ 在 $(1,+\infty)$ 递减，$\phi(x)<\phi(1)=0$，结论成立。
\end{frame}

\begin{frame}{例 2 题目}
\small
已知函数
\[
f(x)=ax^2+(2a-1)x+\ln(x+1),\qquad x>-1,
\]
有两个极值点 $x_1,x_2$。
\begin{enumerate}
  \item 求 $a$ 的取值范围；
  \item 证明：
  \[
  f(x_1)+f(x_2)<2\ln2-\frac54.
  \]
\end{enumerate}
\end{frame}

\begin{frame}[allowframebreaks]{例 2 答案解释}
\small
求导：
\[
f'(x)=2ax+2a-1+\frac1{x+1}
=\frac{2a(x+1)^2-(x+1)+1}{x+1}.
\]
令 $t=x+1>0$，则导函数零点对应
\[
2at^2-t+1=0.
\]
有两个正根需
\[
\Delta=1-8a>0,\qquad a>0.
\]
故
\[
\boxed{0<a<\frac18}.
\]

\framebreak

由韦达定理，
\[
(x_1+1)+(x_2+1)=\frac1{2a},\qquad
(x_1+1)(x_2+1)=\frac1{2a}.
\]
从而
\[
x_1+x_2=\frac1{2a}-2,\qquad x_1x_2=1.
\]
整理可得
\[
f(x_1)+f(x_2)
 =-\frac1{4a}-2a-\ln(2a)+1.
\]
令 $t=2a$，则 $0<t<\frac14$，
\[
F(t)=-\frac1{2t}-t-\ln t+1.
\]
有
\[
F'(t)=\frac1{2t^2}-1-\frac1t
=\frac12\left(\frac1t-1\right)^2-\frac32>0
\]
在 $0<t<\frac14$ 成立，故
\[
F(t)<F\left(\frac14\right)=2\ln2-\frac54.
\]
\end{frame}

\begin{frame}{例 3 题目}
\small
已知函数
\[
f(x)=2x-a\ln x-\frac1x,\qquad x>0,
\]
有两个不同的极值点 $x_1,x_2$，且 $x_1>x_2$。
\begin{enumerate}
  \item 求实数 $a$ 的取值范围；
  \item 若 $a>3$，证明：$x_1>1$，且
  \[
  \frac{f(x_1)-f(x_2)}{x_1+x_2}<\frac43-2\ln2.
  \]
\end{enumerate}
\end{frame}

\begin{frame}[allowframebreaks]{例 3 答案解释}
\small
求导：
\[
f'(x)=2-\frac ax+\frac1{x^2}
=\frac{2x^2-ax+1}{x^2}.
\]
存在两个不同极值点等价于
\[
2x^2-ax+1=0
\]
有两个正根，故
\[
\Delta=a^2-8>0,\qquad a>0.
\]
所以
\[
\boxed{a>2\sqrt2}.
\]

\framebreak

当 $a>3$ 时，两根满足
\[
x_1+x_2=\frac a2,\qquad x_1x_2=\frac12.
\]
且
\[
x_1=\frac{a+\sqrt{a^2-8}}4>1.
\]
计算
\[
\frac{f(x_1)-f(x_2)}{x_1+x_2}
=\frac{4(x_1-x_2)}{x_1+x_2}-2\ln\frac{x_1}{x_2}.
\]
令
\[
t=\frac{x_1}{x_2}>2.
\]
则
\[
\frac{4(x_1-x_2)}{x_1+x_2}-2\ln\frac{x_1}{x_2}
=\frac{4(t-1)}{t+1}-2\ln t.
\]
设
\[
H(t)=\frac{4(t-1)}{t+1}-2\ln t,\qquad t>2.
\]
则
\[
H'(t)=\frac8{(t+1)^2}-\frac2t
=-\frac{2(t-1)^2}{t(t+1)^2}<0.
\]
所以
\[
H(t)<H(2)=\frac43-2\ln2.
\]
\end{frame}

\begin{frame}{例 6 题目}
\small
（对数均值不等式法）仍设
\[
f(x)=\frac1x-x+a\ln x.
\]
\begin{enumerate}
  \item 讨论 $f(x)$ 的单调性；
  \item 若 $f(x)$ 存在两个极值点 $x_1,x_2$，证明：
  \[
  \frac{f(x_1)-f(x_2)}{x_1-x_2}<a-2.
  \]
\end{enumerate}
\ocrnote{原 PDF 例题编号从例 3 跳到例 6，本文保留原编号；本题与例 1 同题不同法。}
\end{frame}

\begin{frame}[allowframebreaks]{例 6 答案解释}
\small
单调性同例 1：
\[
f'(x)=-\frac{x^2-ax+1}{x^2}.
\]
当 $a\le2$ 时 $f$ 递减；当 $a>2$ 时存在两个极值点，且其横坐标 $x_1,x_2$ 满足
\[
x_1x_2=1.
\]

\auditnote{原答案第 (1) 问写“略”，本页补充了单调性讨论。}

\framebreak

当 $a>2$ 时，
\[
\frac{f(x_1)-f(x_2)}{x_1-x_2}
=-\frac1{x_1x_2}+a\frac{\ln x_1-\ln x_2}{x_1-x_2}-1.
\]
由对数均值不等式，
\[
\frac{\ln x_1-\ln x_2}{x_1-x_2}<\frac1{\sqrt{x_1x_2}}
\]
在 $x_1\ne x_2$ 时成立。代入 $x_1x_2=1$ 得
\[
\frac{f(x_1)-f(x_2)}{x_1-x_2}
< -1+a-1=a-2.
\]
\end{frame}

\begin{frame}{例 7 题目}
\small
设函数
\[
f(x)=(ax+1)\ln x-\frac{x^2}{2}-ax+a+\frac12,\qquad a\in\R,\ x>0.
\]
\begin{enumerate}
  \item 当 $a=2$ 时，证明：$f(x)\le0$；
  \item 已知 $f(x)$ 恰好有 $3$ 个极值点 $x_1,x_2,x_3$，且 $x_1<x_2<x_3$。
  \begin{enumerate}
    \item 求实数 $a$ 的取值范围；
    \item 证明：
    \[
    \frac{x_1}{x_2}+\frac{x_2}{x_3}+\frac{x_3}{x_1}>a^2-4a+7.
    \]
  \end{enumerate}
\end{enumerate}
\end{frame}

\begin{frame}[allowframebreaks]{例 7 答案解释}
\small
先求导：
\[
f'(x)=a\ln x+\frac1x-x.
\]
当 $a=2$ 时，
\[
f'(x)=2\ln x+\frac1x-x.
\]
令
\[
\Phi(x)=2\ln x+\frac1x-x.
\]
则
\[
\Phi'(x)=\frac2x-\frac1{x^2}-1=-\frac{(x-1)^2}{x^2}\le0,
\]
且 $\Phi(1)=0$。因此 $\Phi(x)>0$ 对 $0<x<1$ 成立，$\Phi(x)<0$ 对 $x>1$ 成立。
又 $f(1)=0$，故 $f$ 在 $(0,1)$ 上递增，在 $(1,+\infty)$ 上递减，最大值为 $0$，即
\[
f(x)\le0.
\]

\framebreak

若 $f$ 恰有三个极值点，则 $f'(x)=0$ 有三个正根。由于
\[
f'(1)=0,
\]
故 $x_2=1$，另外两根满足
\[
a\ln x+\frac1x-x=0.
\]
记
\[
\Psi(x)=x-\frac1x-a\ln x.
\]
则 $\Psi(1)=0$，且若 $\Psi(x)=0$，则 $\Psi(1/x)=0$，所以另两根互为倒数：
\[
x_1x_3=1,\qquad x_1<1<x_3.
\]
进一步
\[
\Psi'(x)=1+\frac1{x^2}-\frac ax=\frac{x^2-ax+1}{x^2}.
\]
要有除 $1$ 外的两个零点，需 $a>2$。

\framebreak

第二个结论中，因 $x_2=1,\ x_1x_3=1$，
\[
\frac{x_1}{x_2}+\frac{x_2}{x_3}+\frac{x_3}{x_1}
=x_1+\frac1{x_3}+\frac{x_3}{x_1}
=2x_1+x_3^2.
\]
又 $x_1=1/x_3$，可写为
\[
x_3^2+\frac2{x_3}.
\]
先证 $x_3>a-1$。令 $y=a-1>1$，则
\[
\Psi(y)=y-\frac1y-(y+1)\ln y<0,
\]
因为 $\ln y>1-\frac1y$。又 $\Psi(x)\to+\infty$ $(x\to+\infty)$，且 $x_3$ 是 $1$ 右侧的零点，所以 $x_3>a-1$。
令
\[
H(x)=x^2+\frac2x,\qquad x>1.
\]
则
\[
H'(x)=2x-\frac2{x^2}>0.
\]
于是
\[
x_3^2+\frac2{x_3}>(a-1)^2+\frac2{a-1}.
\]
而
\[
(a-1)^2+\frac2{a-1}>a^2-4a+7
\]
等价于
\[
\frac{2(a-2)^2}{a-1}>0,
\]
对 $a>2$ 成立。

\auditnote{原 PDF 此题解析处大量公式缺失，只保留了“故 $a>2$”和目标变形的片段。本版本补全了第 (1) 问和三极值点结构证明。}
\end{frame}

\begin{frame}{例 8 题目}
\small
定义：若函数 $f(x)$ 在定义域内存在极大值 $f(x_1)$ 和极小值 $f(x_2)$，且存在常数 $k$，使
\[
f(x_1)-f(x_2)=k(x_1-x_2),
\]
则称 $f(x)$ 为极值可差比函数，$k$ 称为该函数的极值差比系数。

已知
\[
f(x)=x-\frac1x-a\ln x,\qquad x>0.
\]
\begin{enumerate}
  \item 当 $a=\frac52$ 时，判断 $f(x)$ 是否为极值可差比函数，并说明理由；
  \item 是否存在 $a$ 使 $f(x)$ 的极值差比系数为 $2-a$？若存在求 $a$，若不存在说明理由；
  \item 若 $\frac{3\sqrt2}{2}\le a\le\frac52$，求极值差比系数的取值范围。
\end{enumerate}
\end{frame}

\begin{frame}[allowframebreaks]{例 8 答案解释}
\small
求导：
\[
f'(x)=1+\frac1{x^2}-\frac ax=\frac{x^2-ax+1}{x^2}.
\]
当 $a=\frac52$ 时，
\[
f'(x)=\frac{(2x-1)(x-2)}{2x^2}.
\]
因此 $f$ 在 $(0,\frac12)$ 和 $(2,+\infty)$ 上递增，在 $(\frac12,2)$ 上递减。
极大值点为 $x_1=\frac12$，极小值点为 $x_2=2$，所以
\[
f\left(\frac12\right)-f(2)
=\left(2-\frac{10}3\ln2\right)\left(\frac12-2\right).
\]
故存在极值差比系数，$f$ 是极值可差比函数。

\framebreak

一般地，若有两个极值点，则
\[
x_1+x_2=a,\qquad x_1x_2=1,\qquad a>2.
\]
不妨设 $0<x_1<x_2$。极值差比系数为
\[
k=\frac{f(x_1)-f(x_2)}{x_1-x_2}
=2-\frac{a}{x_1-x_2}\ln\frac{x_1}{x_2}.
\]
若 $k=2-a$，则
\[
\frac1{x_1-x_2}\ln\frac{x_1}{x_2}=1.
\]
代入 $x_1=1/x_2$，等价于
\[
x_2-\frac1{x_2}-2\ln x_2=0,\qquad x_2>1.
\]
令
\[
\phi(x)=x-\frac1x-2\ln x.
\]
则
\[
\phi'(x)=\frac{(x-1)^2}{x^2}>0\quad(x>1),
\]
且 $\phi(1)=0$，故 $x>1$ 时 $\phi(x)>0$，无解。因此不存在这样的 $a$。

\framebreak

设
\[
t=\frac{x_1}{x_2}\in(0,1).
\]
由 $x_1x_2=1$ 得
\[
a^2=(x_1+x_2)^2=t+\frac1t+2.
\]
极值差比系数可写为
\[
p(t)=2-\frac{t+1}{t-1}\ln t.
\]
当
\[
\frac{3\sqrt2}{2}\le a\le\frac52
\]
时，
\[
\frac14\le t\le\frac12.
\]
计算
\[
p'(t)=\frac{2\ln t+\frac1t-t}{(t-1)^2}.
\]
令
\[
q(t)=2\ln t+\frac1t-t.
\]
在 $[\frac14,\frac12]$ 上可证 $q(t)>0$，故 $p(t)$ 递增。
因此
\[
p\left(\frac14\right)\le p(t)\le p\left(\frac12\right),
\]
即
\[
\boxed{2-\frac{10}{3}\ln2\le k\le 2-3\ln2}.
\]
\end{frame}

\subsection{巩固练习}

\begin{frame}{巩固 1 题目}
\small
（安徽省合肥市 2025 届高三一模）已知函数
\[
f(x)=\ln x-a\left(x-\frac1x\right),\qquad a>0,\ x>0.
\]
\begin{enumerate}
  \item 讨论 $f(x)$ 的单调性；
  \item 若函数 $f(x)$ 有两个极值点 $x_1,x_2$，且 $x_1<x_2$，证明：
  \[
  f(x_1)+f(x_2)+f(x_1+x_2)>\ln2-\frac34.
  \]
\end{enumerate}
\ocrnote{原 PDF 题干处有乱码，参考答案中给出了清晰版本，本文按参考答案校正。}
\end{frame}

\begin{frame}[allowframebreaks]{巩固 1 答案解释}
\small
有
\[
f'(x)=\frac1x-a-\frac a{x^2}
=\frac{-ax^2+x-a}{x^2}.
\]
令
\[
P(x)=-ax^2+x-a.
\]
若
\[
1-4a^2\le0\Longleftrightarrow a\ge\frac12,
\]
则 $P(x)\le0$，故 $f$ 在 $(0,+\infty)$ 上单调递减。

若 $0<a<\frac12$，则 $P$ 有两个正根
\[
x_1=\frac{1-\sqrt{1-4a^2}}{2a},\qquad
x_2=\frac{1+\sqrt{1-4a^2}}{2a}.
\]
此时 $f$ 在 $(0,x_1)$、$(x_2,+\infty)$ 上递减，在 $(x_1,x_2)$ 上递增。

\framebreak

若有两个极值点，则 $0<a<\frac12$，且
\[
x_1x_2=1,\qquad x_1+x_2=\frac1a.
\]
于是
\[
f(x_1)+f(x_2)+f(x_1+x_2)
=f(x_1)+f(1/x_1)+f(1/a)
=f(1/a).
\]
计算
\[
f(1/a)=\ln\frac1a-a\left(\frac1a-a\right)
 =-\ln a-1+a^2.
\]
令
\[
h(a)=-\ln a-1+a^2,\qquad 0<a<\frac12.
\]
则
\[
h'(a)=-\frac1a+2a<0.
\]
故 $h$ 在 $(0,\frac12)$ 上递减，
\[
h(a)>h\left(\frac12\right)=\ln2-\frac34.
\]
结论成立。
\end{frame}

\begin{frame}{巩固 2 题目}
\small
已知函数
\[
f(x)=(x-a)\ln x-x\ln a,\qquad a>0,\ x>0.
\]
\begin{enumerate}
  \item 求 $f(x)$ 的极值；
  \item 设
  \[
  g(x)=f(x)+f\left(\frac1x\right)
  \]
  有三个不同的极值点 $x_1,x_2,x_3$。
  \begin{enumerate}
    \item 求实数 $a$ 的取值范围；
    \item 证明：
    \[
    x_1^2+x_2^2+x_3^2>3.
    \]
  \end{enumerate}
\end{enumerate}
\end{frame}

\begin{frame}[allowframebreaks]{巩固 2 答案解释}
\small
求导：
\[
f'(x)=\ln x+\frac{x-a}{x}-\ln a
=\ln x-\frac ax+1-\ln a.
\]
由于 $f'$ 在 $(0,+\infty)$ 上递增，且
\[
f'(a)=0,
\]
故 $f$ 在 $(0,a)$ 上递减，在 $(a,+\infty)$ 上递增，极小值为
\[
\boxed{f(a)=-a\ln a},
\]
无极大值。

\framebreak

令
\[
g(x)=f(x)+f(1/x).
\]
可得
\[
g'(x)=\left(1+\frac1{x^2}\right)\ln x+(\ln a-1)\left(\frac1{x^2}-1\right).
\]
令 $t=x^2>0$，则 $g'(x)=0$ 等价于
\[
h(t)=\ln t-\frac{2(\ln a-1)(t-1)}{t+1}=0.
\]
要有三个不同极值点，$h(t)=0$ 需有三个不同正根。分析
\[
h'(t)=\frac{t^2+(6-4\ln a)t+1}{t(t+1)^2}
\]
可得必须且充分为
\[
\boxed{a>\e^2}.
\]
\auditnote{原答案这里直接给出判别条件，本文补充说明：三极值点来自 $h(t)$ 有三个正根；其导数二次式需有两个正根，并结合 $h(1)=0$ 与端点趋势得到 $a>\e^2$。}

\framebreak

设三个极值点满足 $x_2=1$，并令
\[
t_1=x_1^2,\qquad t_3=x_3^2,\qquad 0<t_1<1<t_3.
\]
由 $h(t_1)=h(t_3)=0$ 得
\[
\varphi(t_1)=\varphi(t_3),\qquad
\varphi(t)=\frac{(t+1)\ln t}{t-1}.
\]
可证 $\varphi$ 在 $(0,1)$ 递减，在 $(1,+\infty)$ 递增。
再令
\[
F(t)=\varphi(t)-\varphi(2-t),\qquad 0<t<1,
\]
可证 $F'(t)<0$，且 $F(1)=0$，故 $F(t)>0$。于是
\[
t_3>2-t_1.
\]
因此
\[
x_1^2+x_2^2+x_3^2=t_1+1+t_3>3.
\]
\end{frame}

\begin{frame}{巩固 3 题目}
\small
定义运算
\[
\begin{vmatrix}a&b\\c&d\end{vmatrix}=ad-bc.
\]
已知函数
\[
f(x)=
\begin{vmatrix}
\ln x&x-1\\
1&a
\end{vmatrix},
\qquad
g(x)=\frac{1-x}{x},\qquad x>0.
\]
\begin{enumerate}
  \item 若函数 $f(x)$ 的最大值为 $0$，求实数 $a$ 的值；
  \item 证明：
  \[
  \left(1+\frac1{2^2}\right)\left(1+\frac1{3^2}\right)\cdots
  \left(1+\frac1{n^2}\right)<\e,\qquad n\ge2;
  \]
  \item 若 $h(x)=f(x)+g(x)$ 存在两个极值点 $x_1,x_2$，证明：
  \[
  \frac{h(x_1)+2x_1-h(x_2)-2x_2}{x_1-x_2}<a.
  \]
\end{enumerate}
\end{frame}

\begin{frame}[allowframebreaks]{巩固 3 答案解释}
\small
由定义
\[
f(x)=a\ln x-x+1.
\]
因此
\[
f'(x)=\frac ax-1.
\]
若 $a\le0$，$f$ 单调递减，不存在最大值。若 $a>0$，$f$ 在 $(0,a)$ 上递增，在 $(a,+\infty)$ 上递减，最大值为
\[
f(a)=a\ln a-a+1.
\]
令
\[
\phi(a)=a\ln a-a+1.
\]
则 $\phi'(a)=\ln a$，$\phi$ 在 $(0,1)$ 上递减，在 $(1,+\infty)$ 上递增，且 $\phi(1)=0$。故
\[
\boxed{a=1}.
\]

\framebreak

由第一问，当 $a=1$ 时
\[
\ln x-x+1\le0,
\]
即
\[
\ln x\le x-1.
\]
当 $k\ge2$ 时，
\[
\ln\left(1+\frac1{k^2}\right)
<\frac1{k^2}
<\frac1{k-1}-\frac1k.
\]
累加得
\[
\sum_{k=2}^n\ln\left(1+\frac1{k^2}\right)
<1-\frac1n<1.
\]
两边取指数：
\[
\prod_{k=2}^n\left(1+\frac1{k^2}\right)<\e.
\]

\framebreak

第三问中
\[
h(x)=f(x)+g(x)=a\ln x-x+\frac1x,
\]
\[
h'(x)=\frac ax-1-\frac1{x^2}
=\frac{-x^2+ax-1}{x^2}.
\]
存在两个极值点等价于
\[
x^2-ax+1=0
\]
有两个不等正根，故
\[
a>2,\qquad x_1x_2=1,\qquad x_1+x_2=a.
\]
计算
\[
\frac{h(x_1)-h(x_2)}{x_1-x_2}
=a\frac{\ln x_1-\ln x_2}{x_1-x_2}-2.
\]
要证
\[
\frac{h(x_1)+2x_1-h(x_2)-2x_2}{x_1-x_2}<a
\]
等价于
\[
\frac{\ln x_1-\ln x_2}{x_1-x_2}<1.
\]
不妨 $0<x_1<1<x_2$，由 $x_1=1/x_2$，只需证
\[
2\ln x_2-x_2+\frac1{x_2}<0.
\]
令
\[
\psi(x)=2\ln x-x+\frac1x,\qquad x>1.
\]
则
\[
\psi'(x)=\frac2x-1-\frac1{x^2}
=-\frac{(x-1)^2}{x^2}<0.
\]
所以 $\psi(x)<\psi(1)=0$，结论成立。
\end{frame}

\begin{frame}{审计汇总}
\small
\begin{tabular}{@{}lll@{}}
\toprule
来源 & 已整理内容 & 备注\\
\midrule
专题一 PDF & 3 个考点，6 道例题，9 道巩固练习 & 巩固 9 与例 6 重复，按原顺序保留\\
专题二 PDF & 基本原理，例 1、2、3、6、7、8，巩固 1、2、3 & 原编号跳号，按原文保留\\
答案审查 & 每题题干页后紧跟答案页 & 有疏漏、乱码、跳步处均用红字标注\\
\bottomrule
\end{tabular}
\end{frame}

\end{document}
